\section{The Bell constraint and the substrate's stance}
\label{sec:bell-constraint}

\subsection{Bell's theorem and its assumptions}

Bell's theorem rules out theories that are simultaneously:
\begin{enumerate}
  \item[(D)] \emph{Deterministic}: hidden variables $\lambda$ fully fix outcomes given the
  measurement settings $a$, $b$.
  \item[(L)] \emph{Local}: the outcome at one wing depends only on the local setting and $\lambda$,
  not on the setting at the distant wing.
  \item[(MI)] \emph{Measurement-independent}: the joint distribution of hidden variables and
  detector settings factorizes, $P(\lambda,a,b)=P(\lambda)P(a,b)$, so the hidden state is
  statistically independent of the measurement choices.
\end{enumerate}
Any theory satisfying all three reproduces only correlations bounded by the Bell inequalities
\citep{Bell1964,CHSH1969}.
The experimental violation of Bell inequalities --- now with all standard loopholes closed
\citep{Hensen2015,Giustina2015,Shalm2015} --- therefore implies that at least one of (D), (L),
(MI) must be given up by any classical-substrate completion.

\subsection{Which assumption the substrate can give up}
\label{subsec:which-assumption}

The substrate satisfies (D) and (L) by construction (\Cref{sec:substrate}):
\begin{itemize}
  \item \textbf{(D) is non-negotiable.} The equations of motion (\Cref{eq:brane-eom}) are
  second-order and deterministic.
  Any probabilistic appearance at the macroscopic level must be emergent
  (\Cref{subsec:born-correlations}), not fundamental (\Cref{subsec:non-probabilistic}).
  Giving up (D) would discard the substrate hypothesis itself.

  \item \textbf{(L) is non-negotiable.} The substrate forces are nearest-neighbor spring contacts.
  Information propagates at finite speed ($c_T$ or $c_L$, \Cref{sec:substrate}).
  Giving up (L) would require adding a new superluminal interaction channel to the substrate ---
  an entirely different model.
\end{itemize}
Therefore (MI) must be relaxed.

\subsection{The probabilistic premise is not fundamental}
\label{subsec:non-probabilistic}

Bell's theorem constrains \emph{probability distributions} over ensembles --- $P(\lambda)$ over
hidden variables, $P(a,b)$ over settings, and the joint $P(\lambda,a,b)$ whose factorization is
assumption (MI). It applies in full to \emph{deterministic} hidden-variable theories, where each
$\lambda$ fixes the outcomes: determinism is one of the cases the theorem targets, not an escape
from it. So a deterministic substrate gains no exemption --- it must still relax one of (D), (L),
(MI). What changes is only \emph{where} the probabilistic description lives.

In this substrate it lives one level up. Its fundamental dynamics are the deterministic field
equations (\Cref{eq:brane-eom}) on a single 4D world-volume: each physical history is one definite
extremal configuration of a local action, with no fundamental randomness, no Born measure, and no
stochastic collapse at the substrate level. The model is therefore \emph{non-probabilistic at its core}, and
in this respect it breaks with the broad family of interpretations --- Copenhagen,
objective-collapse, and any reading that takes the Born rule as a primitive law --- in which
quantum probability is fundamental rather than derived.

Probability re-enters only as emergent, coarse-grained bookkeeping over ensembles of substrate
solutions: the apparent Born-rule statistics of repeated measurements are accounted for as
statistical-mechanical typicality on the carrier amplitude (\Cref{subsec:born-correlations}), not
an input. This does not exempt
the theory from Bell. The emergent statistics must still reproduce the observed correlations, so
the deterministic substrate must still relax one of (D), (L), (MI) at the coarse-grained level ---
which it does, via retrocausality (the subject of the remainder of this section). Being
non-probabilistic underneath does not dodge the theorem; it relocates the probabilistic
description one level up, recasting it as a derived account of ensembles rather than a fundamental
law, and identifies the hidden variable $\lambda$ as an emergent label for a definite substrate
history rather than a primitive random object.

\paragraph{Where the randomness of a photon measurement actually lives.}
This relocation is sharpest for the canonical optical Bell test, in which the entangled carriers
are photons. A free photon is an unconfined linear carrier wave, not a confined soliton, and so
carries no quantization of its own (\Cref{sec:substrate}); the discrete level structure that
quantizes spectra belongs to
the confined matter solitons \citep{MolzbergerMatterMass}. The quantization and the apparent randomness one
associates with light --- a single click, one lump $\hbar\omega$, a definite ``which-detector''
outcome --- are therefore imposed at the detector, where the carrier couples to a bound electron,
not carried by the field in flight. The probabilistic, quantized description attaches to the
emergent matter solitons, while the radiation connecting them remains a definite, continuous
substrate field, and the Born-rule statistics to be recovered (\Cref{subsec:bell-open}) are
located at the detector electron rather than in the light.

\subsection{Why superdeterminism is the wrong way to relax (MI)}
\label{subsec:no-superdeterminism}

One way to relax (MI) is superdeterminism: the hidden state $\lambda$ is correlated with
detector settings $a$, $b$ through a common past boundary condition, so $P(\lambda|a,b)\neq
P(\lambda)$ even without any causal influence between $\lambda$ and the measurement choices.

A first instinct is to rule this out by appeal to forward locality: a nearest-neighbor elastic
medium evolving a Cauchy slice forward has no way to pre-arrange $\lambda$ to anticipate a future
knob setting.
We do \emph{not} use this argument, and the reader should note why: it is inconsistent with the
substrate ontology we actually adopt.
The brane action is time-symmetric and the physical configuration is the stationary point of a
local action over the \emph{entire} 4D world-volume, fixed by boundary data on past \emph{and}
future endcaps (\Cref{sec:worldtube}); it is not a forward-evolved initial-value problem.
An argument that assumes forward evolution would presuppose exactly the causal structure that
retrocausality discards.
The distinction between superdeterminism and retrocausality must therefore be drawn at the level
of \emph{global, all-at-once} dynamics, where both live.

\paragraph{The discriminator is mechanism versus fine-tuning, not global determination.}
It is tempting to think that fixing the whole 4D block at once is already superdeterministic.
It is not. ``All-at-once'' global determination is shared by Hamilton's principle, the classical
electromagnetic action, general relativity, and the Wheeler--Feynman absorber theory --- none of
which is regarded as a conspiracy.
What makes superdeterminism conspiratorial is not that the history is fixed globally, but
\emph{how} the correlation $P(\lambda|a,b)\neq P(\lambda)$ is produced:
\begin{itemize}
  \item \textbf{Superdeterminism} sources the correlation from \emph{special, lawless initial
  data}: a finely arranged configuration on the past boundary, different for every experiment,
  not derived from any dynamical principle, with the setting effectively slaved to the same data
  rather than freely chosen.
  \item \textbf{Retrocausality / all-at-once} sources it from a \emph{uniform local law}: the
  same minimal action functional couples the freely chosen future boundary data to the emission
  state along the worldtube, via local Euler--Lagrange equations.
\end{itemize}
A stationary point of a minimal \emph{local} action is the opposite of a fine-tuned state: it is
the generic, law-selected configuration for whatever boundary data are posed, and the setting
$a$ remains a free boundary input that may be varied at will, with the block re-solving under the
same fixed action.
The substrate's correlation therefore lives in the \emph{law}, not in tuned past data.
This is precisely the ``all-at-once'' (Lagrangian-schema) reading of quantum mechanics developed
by Wharton and collaborators \citep{Wharton2015NotAComputer,WhartonArgaman2020} and the atemporal
notion of lawhood analyzed by Adlam \citep{Adlam2018SpookyTemporal}, both of which dissolve the
superdeterminism charge for global-constraint theories on exactly these grounds.
Superdeterminism, by contrast, hides unexplained correlations in an unobservable initial state
\citep{Hossenfelder2020Superdeterminism} and so undoes the substrate programme's aim of deriving
structure from minimal commitments.

\paragraph{The one residual fine-tuning, named honestly.}
The substrate ontology does retain one genuine tuning: the cosmological boundary data --- the
low-entropy ``past hypothesis'' / Big-Bang-vertex condition that fixes a single direction as time
throughout the brane (\Cref{sec:worldtube}).
This is real, but it is the \emph{universal} initial-condition problem of all of physics, not the
Bell-specific conspiracy.
Superdeterminism requires $\lambda$ to be pre-correlated with each \emph{future detector setting}
in the past data; the all-at-once block does not --- the $\lambda$--$a$ correlation is supplied
by the future boundary through the law.
The substrate thus inherits the ordinary past-hypothesis puzzle while declining the Bell
conspiracy.

\subsection{Retrocausality as the substrate's completion}
\label{subsec:retrocausal-completion}

The substrate relaxes (MI) by letting the future measurement context influence the past state of
the particle along its worldtube:
\begin{quote}
  $P(\lambda | a, b) \neq P(\lambda)$ because the future boundary condition set by
  measurement $a$ back-propagates along the worldtube to the emission event, correlating
  $\lambda$ with $a$ without any spacelike influence.
\end{quote}
This requires only what the brane action already provides: a time-symmetric action
(\Cref{sec:substrate}) whose configurations are fixed by boundary data on \emph{both} temporal
endcaps, so future boundary conditions are as natural as past ones.

\paragraph{Why retrocausality rather than the alternatives.}
The other ways of relaxing the Bell triad are unavailable to this substrate by its ontology, not
by taste. Giving up (D) is excluded by construction. Giving up (L) --- the route of pilot-wave /
Bohmian mechanics and most realist completions --- requires a superluminal channel absent from a
nearest-neighbor elastic lattice (\Cref{subsec:which-assumption}); the Bohmian case additionally
needs a preferred foliation and a configuration-space guidance equation, both alien to the
covariant 4D world-volume. Relaxing (MI) by superdeterminism is rejected as fine-tuned, lawless
initial data rather than a uniform local law (\Cref{subsec:no-superdeterminism}). What the
substrate's own dynamics leaves available is retrocausality, supplied directly by the
time-symmetric action.

\paragraph{The case is constructive, not by elimination.}
We do not rest the argument on having excluded every conceivable alternative. Whether
time-symmetry strictly \emph{forces} retrocausality is contested
\citep{LeiferPusey2017,WoodSpekkens2015}, and a bare elimination argument would inherit that
dispute. The substrate makes a stronger, positive claim instead: it \emph{exhibits} the
retrocausal mechanism and shows that the branching worldtube reproduces the measured
entangled-photon statistics --- no-signalling, the $\cos 2(a-b)$ correlation, the Tsirelson bound,
and the Born weight (\Cref{subsec:born-correlations}). The completion is established by what it
produces, not by what it forbids.

\subsection{Relation to existing retrocausal interpretations}

The retrocausal worldtube reading places itself in the same interpretive family as:
\begin{itemize}
  \item the two-state-vector formalism of Aharonov and collaborators
  \citep{AharonovBergmannLebowitz1964,Aharonov2008TSVF}, which assigns both a forward-propagating
  and a backward-propagating state to each quantum system;
  \item Cramer's transactional interpretation \citep{Cramer1986TI}, in which a transaction is
  completed by a retarded offer wave and an advanced confirmation wave;
  \item the retrocausal models of Price and Wharton \citep{PriceWharton2015Retrocausality}, which
  systematically develop the ontological consequences of time-symmetric dynamics;
  \item Sutherland's time-symmetric Bohmian model \citep{Sutherland2017}, which grounds
  retrocausality in a concrete pilot-wave picture.
\end{itemize}
The present theory differs from all of these in one structural respect: the retrocausal
propagation is not imposed as an interpretive rule on top of a forward-causal dynamics.
It is \emph{built into the substrate action}, which is time-symmetric by construction.
The retrocausal interpretation is therefore not an add-on; it is the direct reading of what the
brane action says.